\documentclass[11pt]{article}

\usepackage[margin=1in]{geometry}
\usepackage{amsmath, amssymb}
\usepackage{xcolor}
\usepackage{listings}
\usepackage{hyperref}
\usepackage{array}
\usepackage{booktabs}

\hypersetup{
    colorlinks=true,
    linkcolor=blue!60!black,
    urlcolor=blue!60!black,
    citecolor=blue!60!black
}

\lstdefinestyle{bugpython}{
  language=Python,
  basicstyle=\ttfamily\small,
  keywordstyle=\color{blue!60!black},
  stringstyle=\color{green!40!black},
  commentstyle=\color{gray!70!black},
  showstringspaces=false,
  breaklines=true,
  frame=single,
  rulecolor=\color{gray!30},
  columns=fullflexible,
  keepspaces=true
}

\title{SymPy Bug Report}
\author{}
\date{}

\begin{document}

\maketitle

\section{Bug 24: solveset returns 1 for asec(x) = 2*pi outside the principal asec range}

\subsection{Minimal Reproducer}

\medskip

\begin{lstlisting}[style=bugpython]
from sympy import Eq, S, asec, pi, simplify, solveset, symbols, N

x = symbols("x")
eq = Eq(asec(x), 2*pi)
sol = solveset(eq, x, domain=S.Complexes)

print("solution:", sol)
print("residual at returned value:", simplify((eq.lhs - eq.rhs).subs(x, 1)))
print("numeric residual:", N((eq.lhs - eq.rhs).subs(x, 1), 50))
\end{lstlisting}

\subsection{Observed Behavior}

\medskip

\begin{lstlisting}[style=bugpython]
solution: {1}
residual at returned value: -2*pi
numeric residual: -6.2831853071795864769252867665590057683943387987502
\end{lstlisting}

SymPy returns the finite solution set \(\{1\}\).  Substituting the returned
value into the original equation does not give zero residual, so the returned
point is not a solution of the equation that was solved.

\subsection{Expected Behavior}

The expected result is \(\varnothing\), represented by SymPy as
\texttt{EmptySet}.  The equation
\[
  \operatorname{asec}(x) = 2\pi
\]
has no solution at \(x=1\), and the returned singleton \(\{1\}\) should not be
reported as the solution set.

\subsection{Mathematical Explanation}

SymPy's principal inverse secant satisfies
\[
  \operatorname{asec}(z) = \arccos\!\left(\frac{1}{z}\right)
\]
on its principal branch.  In particular,
\[
  \operatorname{asec}(1) = \arccos(1) = 0.
\]
Therefore substituting SymPy's returned value gives
\[
  \operatorname{asec}(1) - 2\pi = -2\pi \ne 0.
\]
The incorrect answer is not an equivalent branch representation of the same
solution: \(\operatorname{asec}\) is a principal inverse function, so
\(\operatorname{asec}(1)\) denotes the principal value \(0\), not \(2\pi\).
Applying \(\sec\) to both sides of the equation is therefore not a reversible
operation unless the principal-range condition is preserved or the resulting
candidate is checked against the original equation.

\subsection{Source-Level Diagnosis}

The diagnosis locates the responsible behavior in
\nolinkurl{sympy/solvers/solveset.py}, specifically in the generic inversion
logic of \texttt{\_invert\_complex} around lines 550--557.  The traced call
path is:
\begin{center}
\begin{tabular}{@{}l@{}}
\texttt{\_invert\_complex(asec(\_C) - 2*pi, \{0\})} \\
\(\to\) \texttt{\_invert\_complex(asec(\_C), \{2*pi\})} \\
\(\to\) \texttt{\_invert\_complex(\_C, \{1\})}.
\end{tabular}
\end{center}
According to the diagnosis, this path does not use the branch-aware
\texttt{\_invert\_trig\_hyp\_complex} helper.  Instead, \texttt{\_invert\_complex}
uses the generic one-argument \texttt{.inverse()} path for \(\operatorname{asec}\).
For \(\operatorname{asec}(\_C)\), the inverse method defined in
\nolinkurl{sympy/functions/elementary/trigonometric.py} returns \(\sec\).
The solver consequently rewrites
\[
  \operatorname{asec}(\_C) = 2\pi
  \quad\Longrightarrow\quad
  \_C = \sec(2\pi) = 1,
\]
without checking that \(2\pi\) is in the principal range of
\(\operatorname{asec}\), and without rejecting the finite candidate after
substitution into the original equation.

A plausible fix direction supported by the diagnosis is to avoid treating
inverse trigonometric functions as globally invertible through the generic
\texttt{.inverse()} path.  The implementation could route inverse trigonometric
functions through branch-aware handling, add explicit principal-range checks, or
verify finite candidates against the original equation before returning them.

\begin{lstlisting}[style=bugpython]
# Diagnosed transformation:
# asec(_C) = 2*pi
# _C = asec.inverse()(2*pi) = sec(2*pi) = 1
\end{lstlisting}

\subsection{Additional Failing Instances}

The independently checked family uses targets of the form
\(\operatorname{asec}(x)=2\pi n\).  For every listed nonzero integer \(n\),
SymPy reports \(\{1\}\), but the same principal-value check applies:
\[
  \operatorname{asec}(1) - 2\pi n = -2\pi n \ne 0.
\]
Thus the expected behavior is to avoid returning \(\{1\}\); for these cases,
the independently recorded oracle using
\(\operatorname{asec}(z)=\arccos(1/z)\) also rejects the returned value.

\begin{center}
\footnotesize
\renewcommand{\arraystretch}{1.3}
\begin{tabular}{@{}>{\raggedright\arraybackslash}p{0.35\linewidth}>{\raggedright\arraybackslash}p{0.29\linewidth}>{\raggedright\arraybackslash}p{0.26\linewidth}@{}}
\toprule
\textbf{Input / Expression} & \textbf{SymPy behavior} & \textbf{Expected behavior} \\
\midrule
\(\operatorname{asec}(x)=2\pi\) & returns \(\{1\}\); residual at \(1\) is \(-2\pi\) & must not return \(\{1\}\) \\
\(\operatorname{asec}(x)=4\pi\) & returns \(\{1\}\); residual at \(1\) is \(-4\pi\) & must not return \(\{1\}\) \\
\(\operatorname{asec}(x)=-2\pi\) & returns \(\{1\}\); residual at \(1\) is \(2\pi\) & must not return \(\{1\}\) \\
\(\operatorname{asec}(x)=6\pi\) & returns \(\{1\}\); residual at \(1\) is \(-6\pi\) & must not return \(\{1\}\) \\
\(\operatorname{asec}(x)=8\pi\) & returns \(\{1\}\); residual at \(1\) is \(-8\pi\) & must not return \(\{1\}\) \\
\(\operatorname{asec}(x)=10\pi\) & returns \(\{1\}\); residual at \(1\) is \(-10\pi\) & must not return \(\{1\}\) \\
\bottomrule
\end{tabular}
\end{center}

\subsection{Related Issues Assessment}

The best-effort deduplication check found related public material documenting
the principal behavior of \(\operatorname{asec}\), including that
\(\operatorname{asec}(1)=0\), and general context about branch-sensitive inverse
trigonometric functions.  It did not identify a public SymPy issue, pull
request, or discussion specifically reporting that
\texttt{solveset(Eq(asec(x), 2*pi), x, domain=S.Complexes)} returns
\(\{1\}\).  The deduplication verdict is therefore a new instance of a known
failure mode: global inversion of principal inverse trigonometric functions can
produce extraneous solutions, but this exact minimized \(\operatorname{asec}\)
case is previously unreported in the available evidence and remains
individually actionable as an issue and regression test.

\subsection{Proposed SymPy Regression Test}

\medskip

\begin{lstlisting}[style=bugpython]
import pytest

from sympy import Eq, FiniteSet, S, asec, pi, simplify, solveset, symbols


@pytest.mark.parametrize("n", [1, 2, -1, 3, 4, 5])
def test_solveset_asec_rejects_out_of_range_targets(n):
    x = symbols("x")
    eq = Eq(asec(x), 2*pi*n)
    sol = solveset(eq, x, domain=S.Complexes)

    assert sol != FiniteSet(1)
    assert simplify((eq.lhs - eq.rhs).subs(x, 1)) != 0
\end{lstlisting}

\end{document}
